\begin{abstract}

Large language models show strong general reasoning but struggle with domain-specific technical tasks such as reliability engineering, where answers require precise numerical computation and specialized knowledge. We present a complete pipeline for adapting a general-purpose LLM to the reliability engineering domain, starting from only 56 seed questions extracted from textbooks. We introduce an adapted self-instruct methodology that replaces same-model answer consistency with cross-model verification, combined with paraphrase-based augmentation inspired by MetaMath, to scale the dataset to 866 high-quality question-reasoning-answer triplets. We identify the self-instruct ceiling effect, wherein LLM generators have a tendancy to produce questions systematically easier than real domain problems. We show that paraphrase augmentation is significantly more effective than difficulty scaling for overcoming it. Supervised fine-tuning of Qwen3-8B with QLoRA yields a statistically significant $+18.6\%$ accuracy improvement ($p=0.002$) under 10-fold cross-validation, placing our 8B-parameter model between frontier models o3-mini and Gemini~3.1~Pro on reliability engineering tasks. We further apply Group Relative Policy Optimization (GRPO) to improve reasoning quality beyond what SFT alone achieves. Our results demonstrate that small-seed, LLM-driven data synthesis pipelines can produce effective domain-specific training data, and that combining SFT with reinforcement learning from verifiable rewards is a promising approach for low-resource technical domains.

\TODO{Refine after all results are in.}

\end{abstract}
\renewcommand{\thefootnote}{\fnsymbol{footnote}}
\footnotetext[1]{Code and data: \url{https://github.com/ADnocap/Reliability-Domain-Specific-LLM}}
\footnotetext[2]{Experiments (exp) described in appendix A}
\renewcommand{\thefootnote}{\arabic{footnote}}
